% Suggested placement: physical-validation limits in Online--Offline Equivalence.
% These results concern reverse-candidate sampling, not execution-path agreement.
\subsection{Physical Interpretation of the Reverse-Shock Diagnostics}
\label{subsec:reverse_sampling_limits}

Agreement between online and offline outputs establishes consistency of the
execution paths for the tested configurations. Because both paths use the
same shock-finding calculations, this agreement does not independently
establish the physical accuracy of the reconstructed shock properties.

A separate inspection compared the reverse-shock candidates at $64^3$ and
$128^3$ at $t=0.12$, $0.15$, and $0.20$, with a fixed physical injection
radius of $0.0625$ and the same requested analysis schedule. At $64^3$,
all 48 upstream samples at $t=0.12$ lay inside the injection region, although
the detected surface cells lay outside. No selected endpoints lay inside
that region at the later inspected times or at any inspected time at $128^3$.
Excluding surface cells within the source therefore does not ensure that
the reconstructed states are sampled outside it.

For $\gamma=5/3$, the directly sampled density ratio was compared with
$J_{\rho,\mathrm{RH}}=4\mathcal{M}^2/(\mathcal{M}^2+3)$, using the
pressure-derived Mach number of each sample. The median fractional residual
$J_\rho/J_{\rho,\mathrm{RH}}-1$ was approximately $+49.8\%$, $+36.0\%$,
and $-2.8\%$ at $64^3$, compared with $-5.8\%$, $-7.5\%$, and $-8.7\%$
at $128^3$. These residuals measure internal jump consistency rather than
error against known exact states. The finer run reduced the early discrepancy,
but the later absolute median discrepancy increased. The two resolutions
do not establish convergence. Early reverse-shock thermalization rates and
accumulated energies must therefore be interpreted with these sampling
limitations. Source-overlap flags report this limitation without changing
sample selection or the energy-integration rule.
